\hypertarget{literature-survey}{%
\chapter{Literature Survey}\label{literature-survey}}

\begin{draftnote}

\textbf{Draft status, 03 November 2026.} This prose was written from the
argument and reading notes already recorded in
\texttt{docs/report/chapter2\_plan.md}, which states that every citation
in \texttt{references.bib} was checked against a publisher record or an
author's own copy. \textbf{No factual claim about any paper has been
added here beyond what that plan records.} Two consequences follow, and
both must be dealt with before submission:

\begin{enumerate}
\def\labelenumi{\arabic{enumi}.}
\tightlist
\item
  Nine of the thirteen entries in \texttt{references.bib} still carry a
  TODO on at least one field. Every one is a real publication with a
  confirmed title; each needs a look at the publisher record. Chapter 2
  can be written around them, but the bibliography cannot be printed
  until they are filled.
\item
  The plan's target is 15 to 25 references and the bibliography holds
  13. The seven gaps are named in the plan's own table and restated in
  \S{}2.8 below. Do not pad the count; close the named gaps.
\end{enumerate}

Sections are numbered for a nine-chapter LaTeX report. Citation keys
match \texttt{references.bib}. Replace \texttt{\textbackslash{}cite\{\}}
with the house style if it differs.

\end{draftnote}

\begin{center}\rule{0.5\linewidth}{0.5pt}\end{center}

\hypertarget{what-this-chapter-has-to-establish}{%
\section{What this chapter has to
establish}\label{what-this-chapter-has-to-establish}}

A literature survey in a project report is often a sequence of
summaries. That is not what this chapter is for. Its job is to establish
that the system described in the chapters that follow needed building
--- that the gap it occupies is a real gap in the published work, and
not simply a gap in the author's reading.

The argument runs in four steps, and the chapter is organised around
those steps rather than around chronology. Source-code similarity
detection is a mature field whose best work is explicit that it detects
similarity and not plagiarism. That maturity, however, is confined to
one question: whether two submissions resemble each other. Code
authorship attribution asks a different and more useful question, but
asks it in a form that a university course cannot supply the inputs for.
Detection of machine-generated code exists but depends on training data
that a course cannot supply either. What is missing is not a better
algorithm in any one strand; it is a defensible way of combining them
and reporting the result so that a human being can act on it.

\hypertarget{the-anchoring-result-tools-do-not-find-plagiarism}{%
\section{The anchoring result: tools do not find
plagiarism}\label{the-anchoring-result-tools-do-not-find-plagiarism}}

The most directly relevant published statement of this project's
governing principle is in Novak, Joy and Kermek's systematic review of
source-code similarity detection and the tools used in academia
\cite{novak2019}. Their review concludes that \emph{tools do not find
plagiarism regardless of how good they are}, and that a human must judge
what counts as plagiarism in a given academic context. They identify
JPlag, MOSS and Sherlock Warwick as the most promising of the available
tools, and they report two facts about the field that bear directly on
how this project should be evaluated: most academic tools are never
released, and most are never compared against one another.

This matters more than a design preference being validated. CodeGuard AI
is built as an evidence system rather than a verdict machine: it ranks
submissions for human review and attaches the evidence for its ranking,
and it deliberately does not decide whether plagiarism occurred. Stated
by the author of a project, that is a preference, and it invites the
objection that ranking rather than deciding is a way of avoiding the
hard part of the problem. Stated in \emph{ACM Transactions on Computing
Education} by the authors of a systematic review, it is a position with
support behind it.

It also supplies the justification for the specific output this system
produces. If a tool cannot determine plagiarism, then a single
similarity percentage compared against a threshold is the wrong
interface, because it presents a judgement the tool is not entitled to
make. A ranked score with its evidence attached presents what the tool
actually has: an ordering of where a reviewer's attention is best spent.

\hypertarget{token-and-fingerprint-methods-and-the-question-they-cannot-ask}{%
\section{Token and fingerprint methods, and the question they cannot
ask}\label{token-and-fingerprint-methods-and-the-question-they-cannot-ask}}

The dominant family of deployed tools compares submissions as sequences
of tokens. Schleimer, Wilkerson and Aiken's winnowing algorithm
\cite{schleimer2003} is the fingerprinting scheme behind MOSS. Winnowing
selects the minimum hash from each window of consecutive k-gram hashes,
which yields a guarantee that any shared substring of length at least
\texttt{w\ +\ k\ -\ 1} produces at least one matching fingerprint. That
guarantee is the reason MOSS both works and scales: the index is a small
fraction of the corpus, and the matching property is provable rather
than empirical.

JPlag \cite{prechelt2002} takes a different route to the same question.
It tokenises each submission using a language-aware front end, so that
identifier names and formatting are discarded before comparison, and
then runs greedy string tiling over the resulting token streams to find
maximal matching segments.

Both are strong, mature, and widely used. Both also share a structural
limitation, and it is this limitation that opens the space the present
project occupies: \textbf{they compare submissions against each other.}
They are built to answer the question \emph{did two students hand in the
same work?} Consequently there are three forms of academic misconduct
they cannot see at all:

\begin{itemize}
\tightlist
\item
  work outsourced to a paid third party, where no second copy exists
  anywhere in the cohort;
\item
  work generated by a language model, for the same reason;
\item
  work taken from a previous year's cohort, unless that corpus has been
  retained and is included in the comparison set.
\end{itemize}

A tool that compares submissions pairwise can detect collusion and
nothing else. Collusion is real and is the most common case, which is
why these tools remain valuable. But it is one of the three threats
named in this project's scope, and the other two require a different
question to be asked of a submission: not \emph{does this resemble
another submission?} but \emph{does this resemble its own author?}

\hypertarget{tree-and-ast-methods-layer-1s-ancestry}{%
\section{Tree and AST methods: Layer 1's
ancestry}\label{tree-and-ast-methods-layer-1s-ancestry}}

Moving from token streams to syntax trees is what buys resistance to the
obfuscations a student actually applies. Renaming identifiers,
reordering independent statements and reformatting whitespace all change
the token stream while leaving the abstract syntax tree largely intact,
because an AST discards precisely the surface detail those
transformations alter.

The algorithm this project's structural layer calls directly is APTED,
from Pawlik and Augsten \cite{pawlik2016}. Their earlier RTED paper
\cite{pawlik2015} establishes the dynamic decomposition strategy that
the later work builds on, and the authors ask that both be cited
together; this report does so. DECKARD \cite{jiang2007} takes the
alternative path: rather than computing an edit distance between trees,
it characterises subtrees as numerical count vectors and clusters them
using locality-sensitive hashing, trading exactness for scale. The
distinction is worth stating explicitly in a report whose own Layer 1
chooses exactness and then prefilters to recover scale, which is a
hybrid of the two positions.

Zakeri-Nasrabadi et al.'s systematic review \cite{zakeri2023} supplies
the taxonomy within which Layer 1 can be located: text, token, tree,
graph, metric, learning, image and test-based approaches, with hybrids
the most common category at approximately 28\% of the studies surveyed.
Two of their findings should be carried into this report rather than
discovered under questioning.

First, \textbf{Type IV (semantic) clones remain hard, and many
techniques cannot detect them at all.} Layer 1 shares that limitation.
Two implementations of the same algorithm with different structure will
not be matched by tree edit distance, and the report should say so in
its limitations rather than imply a generality the method does not have.

Second, their observation that close to half of surveyed tools support
Java while many languages have little or no support is the cleanest
available justification for this project measuring both Java and Python
separately instead of asserting language portability. The measurements
in Chapter 7 report the two languages apart from each other for exactly
this reason.

\hypertarget{authorship-attribution-layer-2s-ancestry-and-its-trap}{%
\section{Authorship attribution: Layer 2's ancestry, and its
trap}\label{authorship-attribution-layer-2s-ancestry-and-its-trap}}

Code stylometry asks who wrote a file from its style alone.
Caliskan-Islam et al.~\cite{caliskan2015} is the canonical result,
reporting roughly 94\% accuracy across 1,600 programmers and 98\% across
250, using lexical, layout and AST-derived features on Google Code Jam
submissions. Frantzeskou et al.'s SCAP method \cite{frantzeskou2006}
reaches comparable accuracy from byte-level n-gram author profiles
without parsing the source at all. Abuhamad et al. \cite{abuhamad2019}
extend the task to large author populations using deep feature learning.

This is Layer 2's intellectual ancestry, and it is also the single most
dangerous body of work for this project to be measured against. An
examiner who has read Caliskan-Islam will ask why this project reports a
lift of roughly \textbf{1.4x to 2.5x over chance} where the literature
reports \textbf{94\% accuracy}. The question sounds damning. The answer
is not, and it belongs here in Chapter 2 rather than being offered
defensively in Chapter 7 after the numbers have already been presented.

\textbf{The task is not the same one.} The attribution literature
performs closed-set identification: given a file and a known list of N
candidate authors, select the author. Layer 2 performs anomaly detection
against a single reference: given a file and this student's own earlier
work, is it consistent with that work? There is no candidate list, and
there cannot be one, because the ghostwriter is not enrolled in the
course. The two tasks share features and share nothing else.

\textbf{The feature budget is not the same.} Those systems use hundreds
to thousands of features, including deep AST-derived ones. The scoring
in this project uses four, for two reasons: four is what survived
measurement on real code, and every feature contributing to a flag must
be explainable to a student who contests it. A feature a reviewer cannot
explain is not evidence, whatever it contributes to an accuracy figure.

\textbf{The data is not the same.} Google Code Jam provides many authors
solving \emph{identical} problems, which holds the task constant and
allows style to dominate the signal. Coursework has the same property,
which makes it the fair comparison. The figures this project reports,
however, come from library code across 26 open-source repositories,
where the project's own conventions act as a confound; the honest
within-project figures are the lower ones, and Chapter 7 reports them as
such.

\textbf{The success criterion is not the same.} Attribution is graded on
selecting the correct author. Layer 2 is graded on whether a human
reviewer, shown the evidence, reaches a defensible decision. A 2x lift
that surfaces the right twenty submissions out of five hundred is useful
to a marker. A 94\% attribution accuracy over a candidate set that does
not contain the actual author is not a useful number at all.

Positioning the project against this literature before presenting
results is considerably stronger than defending the gap afterwards, and
this chapter should do so explicitly.

\hypertarget{detection-of-machine-generated-code-and-why-layer-3-was-dropped}{%
\section{Detection of machine-generated code, and why Layer 3 was
dropped}\label{detection-of-machine-generated-code-and-why-layer-3-was-dropped}}

The fourth strand is the newest. Nguyen et al.~\cite{nguyen2023}, with
the journal version appearing as GPTSniffer \cite{gptsniffer2024},
fine-tune CodeBERT to distinguish human-written from ChatGPT-written
code, and report outperforming GPTZero and OpenAI's own text classifier
on that task. Two of their findings are directly relevant: performance
improves when the training data resembles the target code, and it
improves when snippets are classified in pairs rather than individually.

\textbf{The first of those findings is the argument for dropping this
project's third layer, and the report presents it as a decision rather
than as an omission.} A classifier whose accuracy depends on its
training data resembling the target requires a labelled corpus of
student work drawn from the course it will police. No such corpus exists
for this project, and none can be manufactured honestly: generating one
would mean producing synthetic ``AI-written'' student submissions and
treating the result as ground truth.

Without that corpus the layer could still be implemented, but its
performance could never be reported. In a system whose entire premise is
auditable evidence, a layer that produces a number nobody can
characterise is worse than a layer that is absent, because a reviewer
has no basis on which to weigh it. The architecture therefore records
the layer's absence explicitly rather than silently omitting it ---
every detection result stores
\texttt{aiContent:\ \{\ status:\ "skipped",\ reason:\ "Layer\ 3\ is\ not\ implemented\ yet"\ \}}
--- and Chapter 8 lists the layer as future work with the labelled
corpus named as its precondition.

This has a consequence for the score composition that Chapter 4
describes in detail and that is worth flagging here: because the third
layer never runs, the configured weights are renormalised over the
layers that did run, and the headline weights are therefore never the
applied ones.

\hypertarget{the-gap-this-project-occupies}{%
\section{The gap this project
occupies}\label{the-gap-this-project-occupies}}

Drawing the four strands together.

Similarity detection is mature, and the best of it is honest about what
it does: it finds similarity, not plagiarism, and a human decides
\cite{novak2019}. But similarity detection can only see collusion,
because it compares submissions against one another
\cite{schleimer2003, prechelt2002}. Tree-based methods make that
comparison robust to the obfuscations students actually apply
\cite{pawlik2016, jiang2007}, which is a genuine improvement, but it
does not change the question being asked. Authorship attribution asks a
question that could in principle see outsourced work, and asks it well
\cite{caliskan2015, frantzeskou2006, abuhamad2019} --- but it is posed
as closed-set identification against a candidate list, and a course does
not have one. Detection of machine-generated code exists, and its
accuracy depends on in-distribution training data that a course does not
have either \cite{gptsniffer2024}.

The contribution claimed by this project is therefore not a better
algorithm in any one of these strands. It is the combination of
structural similarity with per-student stylistic baselining, composed
into a single ranked and auditable figure with its evidence attached,
measured honestly enough to state where it fails.

That is a defensible contribution for a final-year project, and it is
defensible precisely because it does not claim to beat Caliskan-Islam et
al.~at their own task.

\begin{draftnote}

\textbf{Required revision before submission.} The claim in the paragraph
above is stated as the project's contribution. The measurement reported
in Chapter 7 does not support the stylistic-baselining half of it on the
one corpus built to test it: on 214 spliced pairs the behavioural layer
did not detect partial outsourcing, and the mechanism for that failure
is now understood. \textbf{The wording of this project's claimed novelty
therefore has to be settled with the project guide before this section
is finalised}, and Chapter 7 must be written first. Do not submit \S{}2.7
in its present form.

\end{draftnote}

\hypertarget{limitations-of-this-survey}{%
\section{Limitations of this survey}\label{limitations-of-this-survey}}

Stating these is better than leaving an examiner to find them.

\textbf{Coverage.} The bibliography holds 13 entries against a target of
15 to 25. The gaps are specific rather than general, and each is wanted
for a stated purpose:

\begin{longtable}[]{@{}
  >{\raggedright\arraybackslash}p{(\columnwidth - 2\tabcolsep) * \real{0.5000}}
  >{\raggedright\arraybackslash}p{(\columnwidth - 2\tabcolsep) * \real{0.5000}}@{}}
\toprule\noalign{}
\begin{minipage}[b]{\linewidth}\raggedright
missing
\end{minipage} & \begin{minipage}[b]{\linewidth}\raggedright
wanted for
\end{minipage} \\
\midrule\noalign{}
\endhead
\bottomrule\noalign{}
\endlastfoot
Baxter et al., clone detection using ASTs, ICSM 1998 & the origin of
AST-based clone detection; Layer 1's grandparent \\
Joy and Luck, plagiarism in programming assignments & the
lexical/structural obfuscation taxonomy \cite{novak2019} draws on \\
CCFinder (Kamiya et al.) & the token-based industrial counterpoint to
DECKARD \\
a study of student plagiarism behaviour & Chapter 1 needs evidence that
the problem is real and common, not assumed \\
Sherlock (Warwick) & named in \cite{novak2019}'s top three; worth a
paragraph \\
contract cheating / essay mills & the outsourcing threat model Layer 2
exists to address \\
a reference for one-way ANOVA and the F ratio & Chapter 7's method needs
a citation, not only an explanation \\
\end{longtable}

Baxter et al.~could not be verified at the time of the original search
and \textbf{must not be cited until it has been}.

\textbf{Verification state.} Four entries have every field confirmed
(\texttt{novak2019}, \texttt{prechelt2002}, \texttt{caliskan2015},
\texttt{nguyen2023}). The remaining nine are real publications with
confirmed titles and at least one unconfirmed field. During the first
drafting of the bibliography three fields were written from memory
rather than read from a source --- a DOI on \texttt{schleimer2003} and
the author lists on \texttt{frantzeskou2006} and \texttt{abuhamad2019}
--- and were removed on a second pass. That is recorded because the
bibliography's own header forbids exactly that practice, and the rule
only caught the error because the file was audited after being written.

\textbf{Depth.} This survey positions the project against prior work; it
does not reproduce or re-implement any of it. No claim in this chapter
rests on a comparison this project ran. Where the report compares its
own figures against published ones, as in \S{}2.5, the comparison is of
\emph{tasks} and not of results, and is argued rather than measured.
